\documentclass[11pt,a4paper]{article}
\usepackage[margin=19mm,headheight=15pt]{geometry}
\usepackage{fontspec}
\setmainfont{Noto Serif}
\usepackage{amsmath,amssymb,mathtools}
\usepackage{xcolor}
\usepackage{fancyhdr}
\usepackage{enumitem}
\usepackage{microtype}
\setlength{\parindent}{0pt}
\setlength{\parskip}{6pt}
\definecolor{navy}{HTML}{183153}
\definecolor{lightgray}{HTML}{F3F5F7}
\pagestyle{fancy}\fancyhf{}
\lhead{\small ME 5309 -- Mid-Semester Examination}
\rhead{\small ME 5309}
\cfoot{\thepage}
\begin{document}
\begin{center}
{\Large\bfseries ME 5309 -- Mid-Semester Examination}\\[3pt]
{\large Answer Booklet}\\[8pt]
\end{center}
\textbf{Instructions:} Attempt any three questions. Each question carries 10 marks. Show working where appropriate.
\section*{Question 1 \hfill [10 marks]}
\subsection*{Q1(a) \hfill [4 marks]}

\vspace{2mm}

\subsection*{Q1(b) \hfill [3 marks]}

\[\begin{aligned}Sensitivity = 10 mV/N\end{aligned}\]
The ADC has 12 bits, so there are 4096 levels.

\subsection*{Q1(c) \hfill [3 marks]}

\[\begin{aligned}F = 0.8187, G = 0.0906\end{aligned}\]

\section*{Question 2 \hfill [10 marks]}
\subsection*{Q2(a) \hfill [2 marks]}

Choose x₁ = θ and x₂ = θ̇.
Then ẋ₁ = x₂ and ẋ₂ = [τ - b x₂ - mgℓ sin(x₁)]/J.
At equilibrium:\[\begin{aligned}x_2 = 0 and τ_0 = mg\ell sin \theta_e\end{aligned}\]
Since |sin θ\_e| less than or equal to 1, an equilibrium exists when |τ₀| less than or equal to mgℓ = 2.4525 N m.\par\smallskip At rest $\dot\theta=0$ and $\ddot\theta=0$, so the damping term vanishes. Rearranging gives $\tau_0=mg\ell\sin\theta_e$. Here $mg\ell=1\times9.81\times0.25=2.4525$ N\,m. Since $-1\leq\sin\theta_e\leq1$, the required torque must lie between $-2.4525$ and $+2.4525$ N\,m.

\subsection*{Q2(b) \hfill [2 marks]}

\[\begin{aligned}A = [[0,1], [-mg\ell sin(60^{\circ})/J, -b/J]] = [[0,1], [-106.2, -2.5]],\\ B = [[0],[50]]\end{aligned}\]

\subsection*{Q2(c) \hfill [4 marks]}

At steady state the derivative terms vanish.
The proportional torque is Kp(θd-θ), which must balance mgℓ sin θ.

\subsection*{Q2(d) \hfill [2 marks]}

\vspace{2mm}

\section*{Question 3 \hfill [10 marks]}
\subsection*{Q3(a) \hfill [2 marks]}

\vspace{2mm}

\subsection*{Q3(b) \hfill [2 marks]}

Tustin uses s = (2/T)(z-1)/(z+1).

\subsection*{Q3(c) \hfill [4 marks]}

The characteristic polynomial is s²+(ℓ₁+2)s+(2ℓ₁+ℓ₂).
I match it to s²+10s+100 and get ℓ₁=8, ℓ₂=84.

\subsection*{Q3(d) \hfill [2 marks]}

An observer uses the plant model and known control input to estimate unmeasured states such as velocity, and its poles can be designed.
A simple low-pass filter cannot exploit the state-space model in this way.
A filter may be safer when the model is poorly known or unreliable, because it relies less on model accuracy.\par\smallskip The observer uses the plant model and the measured output to reconstruct the unmeasured state. Its convergence depends on the model and the observer poles. A low-pass filter is simpler, but generally introduces phase lag and does not explicitly use the plant dynamics.

\section*{Question 4 \hfill [10 marks]}
\subsection*{Q4(a) \hfill [2 marks]}

The characteristic polynomial is s²+4s+8, so I report poles -4 and -8.
Unity DC gain requires Nff=8.

\subsection*{Q4(b) \hfill [2 marks]}

The nominal controller has unity DC gain, so θss=1 and the tracking error is zero even after the actuator gain changes.

\subsection*{Q4(c) \hfill [4 marks]}

I obtain s³+(a+bk₂)s²+(c+bk₁)s-bK\_I.
Desired polynomial is s³+8s²+24s+32, so k₂=6, k₁=18, K\_I=-32.

\subsection*{Q4(d) \hfill [2 marks]}

\vspace{2mm}

\end{document}